\clearpage
\SourceLeaf{227}%
\begin{center}\Large\textsc{[Johann Bolyai's Appendix.]}\end{center}
\bigskip
\begin{center}\textsc{Gauss to Gerling. Göttingen, 14 February 1832.}\end{center}
\mbox{.\ .\ .\ .\ .\ .\ .} I should also mention that in the last few days I received a short work from \textls{Hungary} on non-Euclidean geometry in which I find all \textls{my own ideas and results}, developed with great elegance, although presented so concisely that someone unfamiliar with the subject will find it rather difficult to follow. The author is a very young Austrian officer, the son of a friend of my youth with whom I often discussed the subject in 1798, although at that time my ideas were still far from the development and maturity that they have attained through this young man's own reflections. I consider this young geometer v.~\textsc{Bolyai} a genius of the first order. \mbox{.\ .\ .\ .\ .\ .}

\bigskip
\begin{center}\rule{.25\textwidth}{.3pt}\end{center}
\bigskip
\begin{center}\textsc{Gauss to Wolfgang von Bolyai. Göttingen, 6 March 1832.}\end{center}
\mbox{.\ .\ .\ .\ .\ .\ .} Now a few words about your son's work.

If I begin by saying »\textls{that I must not praise it}«, you will no doubt pause for a moment. I cannot do otherwise. To praise it would be to praise myself, for the entire content of the work, the path your
